
\فصل{مقدمه}

شبکه‌های عصبی عمیق \cite{lecun2015deep} امروزه در گستره‌ی وسیعی از کاربردهای واقعی به‌کار گرفته می‌شوند که بسیاری از آن‌ها با داده‌های شخصی و حساس سروکار دارند. هم‌زمان با گسترش این سامانه‌ها، چارچوب‌های قانونی و اخلاقیِ حفاظت از داده نیز سخت‌گیرانه‌تر شده‌اند؛ از جمله‌ی مهم‌ترین آن‌ها می‌توان به {مقررات عمومی حفاظت از داده‌ی اتحادیه‌ی اروپا\پانویس{General Data Protection Regulation (GDPR)}} \cite{gdpr2016} و {قانون حریم خصوصی مصرف‌کننده‌ی کالیفرنیا\پانویس{California Consumer Privacy Act (CCPA)}} \cite{ccpa2018} اشاره کرد. این قوانین برای کاربران {حق فراموش‌شدن\پانویس{Right to be Forgotten}} را به رسمیت می‌شناسند؛ یعنی هر فرد می‌تواند درخواست کند که داده‌هایش از یک سامانه‌ی یادگیری ماشینِ آموزش‌دیده حذف شود. نکته‌ی ظریف اینجاست که صرفِ پاک‌کردن داده از مجموعه‌داده کافی نیست، زیرا مدلی که پیش‌تر روی آن داده آموزش دیده، همچنان ردِّ آن را در پارامترهای خود نگه می‌دارد؛ بنابراین حذف واقعی باید اثرِ داده را از خودِ وزن‌های مدل بزداید.

از سوی دیگر، بسیاری از سامانه‌های هوشمندِ امروزی در محیط‌های {یادگیری پیوسته\پانویس{Continual Learning}} عمل می‌کنند؛ محیط‌هایی که در آن‌ها مدل به‌صورت دنباله‌ای و مرحله‌به‌مرحله با وظایف \پانویس{Task} گوناگون آموزش می‌بیند و داده‌های وظایف پیشین معمولاً دیگر در دسترس نیستند. در رویکرد کلاسیکِ یادگیری پیوسته، چالش اصلی جلوگیری از {فراموشی فاجعه‌بار\پانویس{Catastrophic Forgetting}} \cite{mccloskey1989catastrophic,kirkpatrick2017overcoming} است؛ یعنی جلوگیری از این‌که یادگیری وظیفه‌ی جدید، دانشِ وظایف پیشین را به‌صورت ناخواسته بازنویسی و تخریب کند.

مسئله‌ای که این پژوهش بر آن متمرکز است، چالشی وارون نسبت به فراموشی فاجعه‌بار است و {فراموشی انتخابی\پانویس{Selective Forgetting}} یا {یادگیری‌زدایی ماشینی\پانویس{Machine Unlearning}} \cite{cao2015towards,bourtoule2021machine} نامیده می‌شود. در اینجا هدف، حذفِ عمدی و جراحی‌گونه‌ی دانشِ یک وظیفه‌ی مشخص از مدل است، بدون آن‌که عملکرد مدل روی وظایف باقی‌مانده دچار افتِ شدید شود. به بیان دیگر، فراموشی فاجعه‌بار آسیبی ناخواسته است، اما فراموشی انتخابی پاک‌کردنی خواسته‌شده و کنترل‌شده است که باید آسیب جانبی آن کمینه بماند.

مانعِ اصلی در راهِ فراموشیِ انتخابیِ مؤثر، پدیده‌ی {هم‌پوشانی پارامتری\پانویس{Parameter Overlap}} است. شبکه‌های عصبیِ پرظرفیت، مفاهیم را روی آرایه‌های گسترده و درهم‌تنیده‌ای از پارامترهای مشترک توزیع می‌کنند. در نتیجه، هنگامی که روش‌های متعارفِ یادگیری‌زدایی—مانند {صعود گرادیانی بی‌قید\پانویس{Unconstrained Gradient Ascent}} یا پاک‌سازیِ مبتنی بر اطلاعات فیشر \cite{golatkar2020forgetting}—روی داده‌ی وظیفه‌ی هدف اجرا می‌شوند، به‌طور ناخواسته نمایش‌های کارکردیِ وظایفِ باقی‌مانده‌ای را که به همان پارامترها متکی‌اند نیز مختل می‌کنند و افتِ عملکردِ شدیدی پدید می‌آورند.

چارچوبِ {یادگیری مادام‌العمرِ حریم‌خصوصی‌آگاه\پانویس{Privacy-Aware Lifelong Learning (PALL)}} \cite{ozdenizci2025pall} به‌عنوان سازوکاری برای ماژولارسازیِ دانشِ وظایف معرفی شد. با این حال، پیاده‌سازیِ استانداردِ آن همچنان روی زیرفضاهای بزرگی از شبکه عمل می‌کند و سازوکارِ دقیقی برای کنترلِ هم‌پوشانی‌های سطحِ ویژگی ندارد؛ از همین رو در هر درخواستِ یادگیری‌زدایی، بخشِ قابل‌توجهی از پارامترهای شبکه را به‌روزرسانی می‌کند. این محدودیت، انگیزه‌ی اصلیِ پژوهش حاضر است: طراحیِ سازوکارهایی که هم‌پوشانیِ بحرانی را صریحاً شناسایی و از آن محافظت کنند تا فراموشیِ انتخابی با کمترین آسیب و کمترین هزینه‌ی پارامتری ممکن شود.

\قسمت{تعریف مسئله}

فرض کنید مدلی به‌صورت دنباله‌ای روی مجموعه‌ای از وظایف $\mathcal{T} = \{1, 2, \dots, T\}$ آموزش می‌بیند. هر وظیفه‌ی $t$ دارای مجموعه‌داده‌ی $D_t$ با مجموعه‌ای از کلاس‌های مجزا است، به‌طوری‌که کلاس‌های وظایف مختلف با هم اشتراکی ندارند. در هر گام، سامانه یک {درخواست\پانویس{Request}} دریافت می‌کند که یا از نوع \emph{آموزش} وظیفه‌ی $t$ است و یا از نوع \emph{فراموشی} وظیفه‌ی $t$. هدف، یادگیری وزن‌های شبکه به‌گونه‌ای است که پس از پردازشِ دنباله‌ی درخواست‌ها، عملکرد روی وظایفِ فعالِ باقی‌مانده حفظ شده و دانشِ وظایفِ فراموش‌شده زدوده شود.

در زمان دریافتِ یک درخواستِ فراموشی برای وظیفه‌ی هدف، مجموعه‌ی وظایفِ باقی‌مانده را با $\mathcal{T}_{\text{active}}$ نشان می‌دهیم. دو خواسته‌ی متضاد باید هم‌زمان برآورده شوند:

\begin{itemize}
	\item \textbf{موفقیتِ فراموشی:} دقتِ مدل روی وظیفه‌ی فراموش‌شده باید تا حدِ ممکن به سطحِ تصادفی \پانویس{Chance Level} نزدیک شود؛ یعنی مدل دیگر نباید بتواند کلاس‌های آن وظیفه را تشخیص دهد.
	\item \textbf{حفظِ دانش:} دقتِ مدل روی وظایفِ $\mathcal{T}_{\text{active}}$ باید تا حدِ ممکن دست‌نخورده بماند. بدترین افتِ عملکرد روی این وظایف—که آن را با معیارِ {بیشترین افت\پانویس{WorstDrop}} می‌سنجیم—باید کمینه شود.
\end{itemize}

در نسخه‌ی فعلیِ این پژوهش فرض می‌شود که شناسه‌ی وظیفه \پانویس{Task Identity} در زمان آموزش، فراموشی و ارزیابی در دسترس است؛ یعنی مدل در حالتِ {وظیفه‌آگاه\پانویس{Task-Aware}} کار می‌کند و شناسه‌ی وظیفه‌ای که باید حذف شود از سوی کاربر یا سامانه تعیین می‌گردد. این فرض با بخشی از ادبیاتِ یادگیری پیوسته سازگار است و با سناریوی واقع‌بینانه‌ی حذفِ داده به‌درخواستِ کاربر هم‌خوانی دارد.

به‌طور خلاصه، مسئله‌ی اصلیِ این پژوهش چنین تعریف می‌شود: \emph{چگونه می‌توان دانشِ مربوط به یک وظیفه‌ی مشخص را از مدل حذف کرد، در حالی که پارامترهای مشترک و مهم برای وظایفِ باقی‌مانده تا حدِ ممکن حفظ شوند؟}

\قسمت{چالش‌ها}

پیاده‌سازیِ مؤثرِ فراموشیِ انتخابی در یادگیری پیوسته با مجموعه‌ای از چالش‌های کلیدی روبه‌روست:

\begin{itemize}
	\item \textbf{هم‌پوشانیِ پارامتریِ بحرانی:} برخی پارامترها هم‌زمان برای وظیفه‌ی فراموش‌شونده و برای وظایفِ فعال مهم‌اند. دستکاریِ بی‌محابای این ناحیه، مرزهای تصمیمِ وظایفِ باقی‌مانده را تخریب می‌کند. شناسایی و محافظتِ صریحِ این ناحیه دشوار اما تعیین‌کننده است.

	\item \textbf{حفظِ تعادل میان فراموشی و نگه‌داری:} اگر فراموشی بیش‌ازحد محافظه‌کارانه باشد، دانشِ وظیفه‌ی هدف به‌طور کامل زدوده نمی‌شود؛ و اگر بیش‌ازحد تهاجمی باشد، به وظایفِ فعال آسیب می‌زند. یافتنِ نقطه‌ی تعادل، چالشی اساسی است.

	\item \textbf{ملاحظاتِ حریم خصوصی:} فراموشی باید چنان باشد که نتوان از روی مدل دریافت آیا داده‌ی وظیفه‌ی فراموش‌شده زمانی در آموزش حضور داشته است یا نه. حمله‌هایی مانند {استنتاج عضویت\پانویس{Membership Inference Attack}} و {وارون‌سازی مدل\پانویس{Model Inversion}} این موضوع را می‌سنجند.

	\item \textbf{هزینه‌ی پارامتری و محاسباتی:} روش‌های متعارف در هر درخواستِ فراموشی بخشِ بزرگی از شبکه را به‌روزرسانی می‌کنند که هزینه‌بر است. کاهشِ نسبتِ پارامترهای به‌روزرسانی‌شده، به‌ویژه برای استقرار در سامانه‌های با منابعِ محدود، اهمیت دارد.

	\item \textbf{فقدانِ تضمینِ صوری:} اغلب روش‌های عملی به {فراموشیِ تقریبی\پانویس{Approximate Unlearning}} می‌رسند، نه {فراموشیِ دقیق\پانویس{Exact Unlearning}}؛ بنابراین ارائه‌ی شواهدِ تجربی یا کرانِ نظری برای کیفیتِ فراموشی ضروری است.
\end{itemize}

\قسمت{مشارکت‌های پژوهش}

برای رفعِ محدودیت‌های بالا، در این پژوهش دو روشِ مکمل و یک سازوکارِ نظری ارائه می‌شود:

\begin{enumerate}
	\item \textbf{روشِ \lr{PALL-Modified}:} توسعه‌ای که {هم‌پوشانیِ بحرانی\پانویس{Critical Overlap}} میان پارامترهای فراموش‌شونده و فعال را به‌کمکِ بزرگیِ گرادیان روی حافظه‌ی بازپخش \پانویس{Replay Buffer} صریحاً محاسبه می‌کند و فراموشی را با یک جریمه‌ی تنظیمِ محافظتی روی همان ناحیه انجام می‌دهد. این روش بالاترین دقتِ نگه‌داری را به‌دست می‌آورد.

	\item \textbf{روشِ \lr{PALL-Adapter}:} معماریِ کارآمدِ پارامتری که بدنه‌ی استخراج ویژگی را منجمد می‌کند و یادگیری و فراموشی را از طریقِ {آداپترهای گلوگاهیِ مخصوصِ وظیفه\پانویس{Task-Specific Bottleneck Adapters}} و یک {آداپترِ گلوگاهیِ مشترک\پانویس{Shared Bottleneck Adapter}} هدایت می‌کند و در نتیجه نسبتِ پارامترهای به‌روزرسانی‌شده را به‌شدت کاهش می‌دهد.

	\item \textbf{سازوکارِ سه‌فازیِ ماسکِ نرم:} یک الگوریتمِ بهینه‌سازیِ ماسک‌دار (ریشه‌کنی، کمی‌سازی، و مقیاس‌دهیِ نرم) که به‌روزرسانیِ گرادیان را در حین فراموشی کنترل می‌کند تا ویژگی‌های مشترکِ بحرانی محافظت شوند، به‌همراه تحلیلی نظری که نشان می‌دهد این سازوکار، بیشترین افت روی وظایفِ فعال را به‌طور کران‌دار کاهش می‌دهد.
\end{enumerate}

\قسمت{دستاوردهای پژوهش}

با توجه به چالش‌های مطرح‌شده، هدفِ اصلیِ این پژوهش طراحیِ سازوکارهایی برای فراموشیِ انتخابی است که هم‌زمان دانشِ وظیفه‌ی هدف را به‌طور مؤثر بزدایند و آسیب به وظایفِ باقی‌مانده را کمینه نگه دارند. در این راستا، روشِ نخست یعنی \lr{PALL-Modified} نشان می‌دهد که با محاسبه‌ی صریحِ هم‌پوشانیِ بحرانی و اعمالِ یک جریمه‌ی تنظیمِ محافظتی تنها روی همان ناحیه، می‌توان بالاترین دقتِ نگه‌داری را میان روش‌های مقایسه‌شده به‌دست آورد؛ بدونِ آن‌که فرآیندِ فراموشی، فضای کارکردیِ مشترکِ وظایفِ فعال را تخریب کند.

دستاوردِ کلیدیِ دوم، نشان‌دادنِ این نکته است که می‌توان هزینه‌ی پارامتریِ فراموشی را به‌شدت کاهش داد بدونِ آن‌که پایداری از دست برود. روشِ \lr{PALL-Adapter} با منجمدکردنِ بدنه‌ی شبکه و محدودکردنِ یادگیری و فراموشی به آداپترهای کوچک، نسبتِ پارامترهای به‌روزرسانی‌شده را تا حدودِ ۸۷٪ نسبت به نسخه‌های تمام‌شبکه کاهش می‌دهد و هم‌زمان بیشترین افتِ عملکرد روی وظایفِ نگه‌داری‌شده را در تنظیماتِ ارزیابی‌شده زیرِ یک درصد نگه می‌دارد. این روش نشان می‌دهد که مصالحه‌ی میانِ دقت و کارایی را می‌توان آگاهانه به‌سمتِ کم‌هزینگی و پایداری هدایت کرد؛ ویژگی‌ای که برای استقرارِ مکررِ سامانه‌های حریم‌خصوصی‌محور اهمیت دارد.

سومین دستاورد، پشتوانه‌ی نظریِ سازوکارِ پیشنهادی است. نشان داده می‌شود که ضرب‌گرِ ماسکِ نرم صرفاً یک ترفندِ مهارِ شهودی نیست، بلکه به‌طور کران‌دار بیشترین افتِ عملکرد روی وظایفِ فعال را نسبت به فراموشیِ بی‌قید کاهش می‌دهد. این ادعاها روی محک‌های \lr{CIFAR-10}، \lr{CIFAR-100} و \lr{TinyImageNet} و در مقایسه با مجموعه‌ای از روش‌های پایه به‌صورت تجربی اعتبارسنجی می‌شوند، و رابطه‌ی میانِ میزانِ هم‌پوشانیِ بحرانی و آسیبِ ناشی از فراموشی نیز مورد بررسی قرار می‌گیرد.

\قسمت{ساختار پایان‌نامه}

در فصلِ نخستِ این پایان‌نامه، کلیاتِ پژوهش بیان شد. در این فصل، ضمن معرفیِ اهمیتِ شبکه‌های عصبیِ عمیق در کاربردهای حساس به حریم خصوصی و طرحِ حقِ فراموش‌شدن، مسئله‌ی فراموشیِ انتخابی در یادگیری پیوسته به‌عنوان چالشی وارون نسبت به فراموشیِ فاجعه‌بار معرفی گردید. سپس مانعِ اصلیِ آن، یعنی هم‌پوشانیِ پارامتری، تبیین شد و تعریفِ صوریِ مسئله، چالش‌ها، و دستاوردهای پژوهش ارائه شدند.

فصلِ دوم به مرورِ کارهای پیشین اختصاص دارد. در این فصل، ابتدا پارادایمِ یادگیری پیوسته و روش‌های مبتنی بر بازپخشِ حافظه بررسی می‌شوند؛ سپس راهبردهای گوناگونِ یادگیری‌زدایی ماشینی، از فراموشیِ دقیق تا فراموشیِ تقریبی، تحلیل می‌گردند؛ و در پایان، روش‌های یادگیریِ کارآمدِ پارامتری معرفی می‌شوند تا جایگاهِ پژوهشِ حاضر در میانِ این سه محور روشن شود.

در فصلِ سوم، روشِ پیشنهادی به‌تفصیل معرفی می‌شود. ابتدا محدودیتِ فراموشیِ بی‌قید تشریح می‌گردد، سپس روشِ \lr{PALL-Modified} و سازوکارِ محاسبه و محافظت از هم‌پوشانیِ بحرانی شرح داده می‌شود. در ادامه، معماریِ \lr{PALL-Adapter} و فرآیندِ سه‌فازیِ ماسکِ نرم (ریشه‌کنی، کمی‌سازی، و مقیاس‌دهیِ نرم) ارائه می‌شود و در نهایت تحلیلِ نظریِ کرانِ بیشترین افت آورده می‌شود.

فصلِ چهارم به آزمایش‌ها و نتایج می‌پردازد. در این فصل، مجموعه‌داده‌ها و پروتکلِ آزمایش، معیارهای ارزیابی (شاملِ دقتِ نهایی، فراموشیِ میانگین، بیشترین افت، دقتِ وظیفه‌ی فراموش‌شده، و نسبتِ پارامترهای به‌روزرسانی‌شده) معرفی می‌شوند و سپس نتایجِ تجربیِ روش‌های پیشنهادی در مقایسه با روش‌های پایه گزارش و تحلیل می‌گردد.

در نهایت، فصلِ پنجم به جمع‌بندیِ پژوهش اختصاص دارد. در این فصل، دستاوردهای اصلی مرور می‌شوند، محدودیت‌های روش‌های پیشنهادی—از جمله ماهیتِ تقریبیِ فراموشی—بیان می‌گردد، و مسیرهای پژوهشیِ آینده، مانندِ افزودنِ ارزیابیِ حمله‌ی استنتاج عضویت و حرکت به‌سمتِ تضمین‌های صوریِ حریم خصوصی، مطرح می‌شوند.
